\section{Profiles and angular tests}\label{test:section}

We construct directional tests from normalized regular classes and use
the cap estimates of Proposition~\ref{prep:starting-pairs} to bound their
failure mass. A parallel test counts all occupied class cube centres
satisfying its stated ball and tube conditions; a projection test uses
the full conditional probability in a specified class cube. The
exceptional-direction estimate applies to a spherical subprobability.
At its outer use, this is an unnormalized directional fiber of the
original opposite starting law, restricted to the cap-controlled pair set.

For each frequency shell, a specified product of the masks constructed
below will weight the original product measure restricted to $\Gamma_N$.
The smaller cap-controlled subset is used only to estimate the original
pair mass removed by these masks. The product partition of $\Gamma_N$
permits the shell estimate to use separate spatial restrictions at the
two endpoints. Internal positive counts restrict products of normalized
class laws by the good tests, while internal Fourier energies use endpoint
masks and bounded spatial multipliers.

Let $S=d/2$ and $R=2^N$. We use the separated probabilities from
Proposition~\ref{prep:starting-pairs}, of Frostman exponent $S+\chi$.
The parameters are chosen in the following order:
\[
0<\beta<\min(\chi/10,1/10),\qquad
0<\kappa\ll_d\beta/d,\qquad
0<\sigma\ll_{d,\beta,\kappa}\kappa.
\]
The spatial cutoff parameter $\omega$ in
Proposition~\ref{prep:starting-pairs} is chosen afterwards, as small as
required by these fixed parameters. All estimates concern
sufficiently large $N$ after these choices. The constants multiplying
$\kappa$, $\sigma$, or $\sigma/\kappa$ in exponents are dimensional;
ordinary multiplicative constants and the rates in $o(1)$ may depend
on the fixed parameters and the starting geometry. Use the preparation
Proposition with angular loss $R^\sigma$. Let $\rho$ range over the
normalized classes supplied by applying Lemma~\ref{reg:regularization}
with parameter $\sigma$ to each starting probability, and write $m$ for
its profile. All definitions in this Section are made separately for each
class.

\subsection{Profile increments and conditional laws}

For integer depths set
\begin{equation}\label{test:profile-definitions}
\begin{aligned}
M(n)&=R^{-m(n/N)},& A_n&=m(n/N)-Sn/N,\\
I(g,p)&=m(p/N)-m(g/N),&
k(g,p)&=A_p-\min_{g\le n\le p}A_n.
\end{aligned}
\end{equation}
Every minimum over depths is over integers unless stated otherwise.
The quantity $M(n)$ is the model mass of a depth-$n$ cell, and
$I(g,p)$ records the profile increment between two depths. The residual
$A_n$ subtracts the critical slope $S$, while $k(g,p)$ is the rise of
$A_p$ above its minimum since $g$. These quantities will give the
cancellations in the two directional estimates below.

\begin{lemma}[Profile bounds]\label{test:profile-bounds}
There is a dimensional constant $C_d$ such that
\[
A_0=0,\qquad |A_i-A_j|\le d|i-j|/N,\qquad
A_n\ge\chi n/N-C_d\sigma.
\]
Every occupied depth-$n$ cube has mass
$M(n)R^{\pm C_d\sigma}$.  In particular, a ball of radius
$C2^{-g}$ contains at most
\begin{equation}\label{test:centre-count}
C'R^{I(g,p)+C_d\sigma}
\end{equation}
occupied depth-$p$ cube centres, for $0\le g\le p\le N$.
Here $C$ is any fixed constant and $C'$ may depend on $C$.
For integers $0\le p\le t\le N$ and an occupied depth-$p$ cube $P$,
the full conditional probability $\rho_P=\rho|_P/\rho(P)$,
rescaled by $z\mapsto2^p(z-c_P)$, satisfies the $S$-ball bound down to
$H^{-1}$, where $H=2^{t-p}$, with loss
\begin{equation}\label{test:conditional-ball-loss}
K=R^{A_p-\min_{p\le n\le t}A_n+C_d\sigma}.
\end{equation}
\end{lemma}

\begin{proof}
The first two assertions follow from $m(0)=0$ and the fact that the
slopes of $m$ lie in $[0,d]$.  The lower bound follows from
Equation~\eqref{reg:profile-lower-bound}; the fixed Frostman constant
is absorbed into $R^{C_d\sigma}$ for sufficiently large $N$.
The occupied-cell mass assertion is
Equation~\eqref{reg:cell-mass}.  A ball of radius $C2^{-g}$ meets
$O_{d,C}(1)$ depth-$g$ cubes.  Their full masses are at most
$M(g)R^{C_d\sigma}$ each, whereas each occupied depth-$p$ cube has
mass at least $M(p)R^{-C_d\sigma}$.  Summation gives
Equation~\eqref{test:centre-count}, with a larger $C_d$.

For the last assertion, if $p\le n\le t$, a physical ball of radius
$2^{-n}$ meets a bounded number of depth-$n$ cubes, and division by
$\rho(P)$ gives an upper bound
\[
R^{-[m(n/N)-m(p/N)]+C_d\sigma}
=2^{-S(n-p)}R^{A_p-A_n+C_d\sigma}.
\]
Dyadic comparison of radii yields
Equation~\eqref{test:conditional-ball-loss}.  Only depths through $t$
are used, so no mass estimate below $H^{-1}$ is involved.
\end{proof}

The two mass normalizations will remain distinct: $M(n)$ is within the
stated comparison factor of the actual mass of an occupied depth-$n$ cell,
whereas the probability $\rho_P$ is divided by the actual mass $\rho(P)$.

\subsection{Directional tests}

A descriptor $\mathfrak t$ consists of a test type and a pair of integer
depths. A parallel descriptor $[g,p]$, with $0\le g<p\le N$, has
\emph{anchor} $p$ and \emph{nominal length} $p-g$. A projection descriptor
$[p,t]$, with $0\le p<t\le N$, has anchor $p$ and nominal length $t-p$.
The type is part of the identity even when the displayed depth pairs
coincide. These descriptors specify individual tests; a later construction
will select fixed lists of them for depth intervals.

Set
\[
J_0=\left\lceil100(1/\kappa+1/\sigma)+100\right\rceil,
\qquad
W_i(n)=(2-i/J_0)2^{-n}R^{500\sigma},
\quad0\le i\le J_0.
\]
$J_0$ is the maximum test index, fixed independently of $N$ once
$\kappa,\sigma$ are fixed. Thus there are $J_0+1$ good-test levels;
the masks below are defined for $i<J_0$.

\begin{definition}[Good directions]\label{test:good-directions}
Fix a class, a descriptor $\mathfrak t$ with anchor $p$, and an occupied
depth-$p$ cube $P$, of centre $c_P$.  Its level-$i$ good set
$G_{i,\mathfrak t}(P)\subset S^{d-1}$ is defined as follows, with one
sufficiently large dimensional constant $C_d$ for all tests.
\begin{enumerate}
\item For a parallel descriptor $[g,p]$, a direction $w$ is good if
      the number of occupied depth-$p$ cube centres in
      $B(c_P,10(d+1)2^{-g})$ whose distance to $c_P+\R w$ is at most
      $W_i(p)$ is at most
      \begin{equation}\label{test:parallel-threshold}
      R^{k(g,p)+C_d\sigma}.
      \end{equation}
\item For a projection descriptor $[p,t]$, let
      $\rho_P=\rho|_P/\rho(P)$.  A direction $w$ is good if
      \begin{equation}\label{test:projection-threshold}
      (\rho_P\otimes\rho_P)
      \{|(z-z')\cdot w|\le W_i(t)\}
      \le2^{-(t-p)}
           R^{A_p-\min_{p\le n\le t}A_n+C_d\sigma}.
      \end{equation}
\end{enumerate}
For an unoccupied anchor cube set $G_{i,\mathfrak t}(P)=\varnothing$.
For $x$ in the root cube, write
$G_{i,\mathfrak t}(x)=G_{i,\mathfrak t}(P)$ for its unique depth-$p$
descendant $P$, and set $G_{i,\mathfrak t}(x)=\varnothing$ outside the root.
The parallel test uses all occupied class centres in its stated ball and
tube. The projection test uses the full conditional class probability
in $P$. Subsequent pair
restrictions do not change either test.
\end{definition}

\subsection{A Fourier estimate for half-dimensional directions}

The two ball bounds below use the critical exponent $S=d/2$.
Their exponents sum to $d$, so the annular estimate contributes a
constant at each dyadic frequency scale.

\begin{lemma}[Half-dimensional Fourier estimate]\label{test:half-dimensional-fourier}
Let $1\le H\le R$, $K,L\ge1$, and $S=d/2$.  Suppose $\alpha$ is a
probability in a fixed bounded subset of $\R^d$ with
\[
\alpha(B(x,r))\le Kr^S\qquad(H^{-1}\le r\le1),
\]
and $\xi$ is a subprobability on $S^{d-1}$ with
\[
\xi(B(w,r))\le Lr^S\qquad(R^{-1}\le r\le1).
\]
Then, with Fourier convention
$\widehat\alpha(v)=\int e^{-iz\cdot v}\,d\alpha(z)$,
\begin{equation}\label{test:fourier-estimate}
\int_{-H}^{H}\int_{S^{d-1}}|\widehat\alpha(sw)|^2\,d\xi(w)\,ds
 \le C LK(1+\log H).
\end{equation}
The constant depends only on $d$ and the fixed bound on the support's
diameter.  In particular it does not require ball bounds below
$H^{-1}$ or any Fourier decay.
\end{lemma}

\begin{proof}
The interval $|s|\le1$ costs at most $2$.  For a dyadic $B\ge1$,
let $\lambda_B$ be the pushforward of
\[
\one_{\{B\le|s|\le\min(2B,H)\}}\,ds\,d\xi(w)
\]
under $(s,w)\mapsto sw$.  It is supported in the annulus
$B\le|v|\le2B$ and has total mass at most $2B$.
The preimage of a unit frequency ball has radial length $O(1)$ and
lies in at most two spherical caps of radius $O(B^{-1})$.
Since $B\le H\le R$, the cap bound applies.  A fixed covering handles
radii differing from $B^{-1}$ by fixed constants, and the case
$B=O(1)$ uses the total mass bound.  Thus
\begin{equation}\label{test:frequency-unit-ball}
\sup_v\lambda_B(B(v,1))\le C LB^{-S}.
\end{equation}

Translation of $\alpha$ changes only the phase of its transform, so
assume its support is in a fixed ball.  Choose a smooth compactly
supported spatial function equal to one on that ball.  Fourier
convolution and the weighted Cauchy--Schwarz inequality give a
nonnegative, integrable, rapidly decreasing kernel $\Phi$, independent
of $B$, for which
\[
|\widehat\alpha(v)|^2
 \le C\int_{\R^d}|\widehat\alpha(u)|^2\Phi(v-u)\,du.
\]
Indeed one may take a rapidly decreasing majorant for the absolute
value of the cutoff's Fourier transform.  Integrating against
$\lambda_B$ and writing $W_B=\Phi*\lambda_B$, we obtain
\[
\int|\widehat\alpha(v)|^2\,d\lambda_B(v)
 \le C\int|\widehat\alpha(u)|^2W_B(u)\,du.
\]
Summing Equation~\eqref{test:frequency-unit-ball} over unit cubes
against the decay of $\Phi$ gives $W_B\le CLB^{-S}$.  On $|u|>4B$,
the support and total mass of $\lambda_B$ give, for every fixed $M$,
\begin{equation}\label{test:frequency-tail}
W_B(u)\le C_M B(1+|u|)^{-M}.
\end{equation}

A Gaussian majorant and its positive Fourier transform show that
\[
\begin{split}
\int_{|u|\le4B}|\widehat\alpha(u)|^2\,du
&\le CB^d\iint e^{-cB^2|z-z'|^2}\,d\alpha(z)\,d\alpha(z')\\
&\le CKB^{d-S}.
\end{split}
\]
For completeness, the last bound follows by splitting the $z'$ integral
into the ball of radius $B^{-1}$ and annuli with radii
$2^jB^{-1}$.  Every radius is at least $H^{-1}$.  Up to radius one
use the assumed ball bound; above radius one use total mass at most
one, which is bounded by $Kr^S$.  The resulting sum is bounded by
$CKB^{-S}\sum_{j\ge0}2^{jS}e^{-c'4^j}$.
The part with $|u|\le4B$ therefore costs at most
\[
CLK B^{d-2S}=CLK.
\]
For the complement, use $|\widehat\alpha|\le1$ and
Equation~\eqref{test:frequency-tail}, choosing $M>d+1$.  Its
contribution is at most $C B^{d+1-M}\le C$, which is absorbed by
$CLK$.  Summing over the $O(1+\log H)$ dyadic shells proves
Equation~\eqref{test:fourier-estimate}.
\end{proof}

\subsection{Exceptional mass of the tests}

\begin{lemma}[Exceptional directions]\label{test:exceptional-directions}
The constant $C_d$ in Definition~\ref{test:good-directions} can be
chosen depending only on $d$ so that the following holds for all fixed
parameters and sufficiently large $N$.  If $\xi$ is any spherical
subprobability satisfying
\[
\xi(B(w,r))\le R^\sigma r^S
       \qquad(R^{-1}\le r\le1),
\]
then every level-zero test, uniformly in the class and the occupied
anchor cube, has
\begin{equation}\label{test:exceptional-bound}
\xi(S^{d-1}\setminus G_{0,\mathfrak t}(P))\le R^{-\sigma}.
\end{equation}
\end{lemma}

\begin{proof}
First take a parallel descriptor.  The anchor centre itself costs at
most one in the count.  Divide the other centres into annuli of
distance comparable to $2^{-h}$ from $c_P$, where
$g-O_d(1)\le h\le p+O_d(1)$.  A given centre in this annulus is counted
only for directions in two caps of radius
\[
O(2^{-(p-h)}R^{500\sigma}).
\]
For radii at most one, the angular mass is at most
$C R^{(1+500S)\sigma}2^{-S(p-h)}$.  For larger radii the same bound,
after adjusting the constant, follows from the subprobability
property.  All small radii used are at least a fixed multiple of
$R^{-1}$; changing this multiple costs only a fixed cap covering.
The number of centres is bounded by
$C R^{I(h,p)+C_d'\sigma}$ by
Equation~\eqref{test:centre-count}.  If an endpoint annulus has
$h$ outside $[g,p]$, replace it by the nearer endpoint, which changes
scales only by a dimensional constant.  The identity
\[
I(h,p)-S(p-h)/N=A_p-A_h\le k(g,p)
\]
then gives expected count at most
\[
C(N+1)R^{k(g,p)+C_d''\sigma}.
\]
Here $C_d''$ is dimensional, including the factor $1+500S$.
Choose $C_d>C_d''+3$.  For fixed $\sigma$ and sufficiently large $N$,
Markov's inequality at threshold
Equation~\eqref{test:parallel-threshold} gives
Equation~\eqref{test:exceptional-bound}.

For a projection descriptor, rescale $P$ to unit size and let
$\alpha$ denote its conditional law.  Set $H=2^{t-p}$ and
$\delta_A=A_p-\min_{p\le n\le t}A_n\ge0$.  By
Lemma~\ref{test:profile-bounds}, its $S$-ball loss is
$K=R^{\delta_A+C_d'\sigma}$.  The rescaled tolerance of the
level-zero test is
$\varepsilon=2H^{-1}R^{500\sigma}$.  A constant multiple of
\[
F_\varepsilon(u)=
 \left(\frac{\sin(u/(2\varepsilon))}{u/(2\varepsilon)}\right)^2
\]
majorizes $\one_{[-\varepsilon,\varepsilon]}$.  Its Fourier transform
is nonnegative, supported in
$[-\varepsilon^{-1},\varepsilon^{-1}]\subseteq[-H,H]$, and bounded
by $C\varepsilon$.  Fourier inversion and Tonelli's Theorem thus
bound the mean collision mass by
\[
\begin{split}
\int (\alpha\otimes\alpha)
 \{|(z-z')\cdot w|\le\varepsilon\}\,d\xi(w)
&\le C\varepsilon
 \int_{-H}^H\int|\widehat\alpha(sw)|^2\,d\xi(w)\,ds\\
&\le C H^{-1}(1+\log H)
 R^{\delta_A+(501+C_d')\sigma}.
\end{split}
\]
The last inequality is Lemma~\ref{test:half-dimensional-fourier}
with $L=R^\sigma$.  Increasing the same dimensional $C_d$ so that
$C_d>504+C_d'$ absorbs constants and $1+\log H$ and gives
Equation~\eqref{test:exceptional-bound} by Markov's inequality.
This argument also covers large tolerances; alternatively, if the
threshold is at least one, the test holds for every direction.
The choices of $C_d$ depend only on previously fixed dimensional
constants, not on any small parameter.
\end{proof}

\subsection{Smooth masks and products}

The width decreases as the index increases, while the count or mass
threshold stays fixed.  Consequently the sets $G_{i,\mathfrak t}$ increase with
$i$.  The considerable gap between the width exponent $500\sigma$
and the smoothing exponent $50\sigma$ gives a uniform buffer.

\begin{lemma}[Masks with a directional buffer]\label{test:smooth-masks}
For a descriptor $\mathfrak t$ of nominal length $l'$, put
\[
D_{l'}=\max(1,2^{l'}R^{-50\sigma}).
\]
For all sufficiently large $N$, every $0\le i<J_0$ admits a mask
$w_{i,\mathfrak t}(x,w)\in[0,1]$, measurable in $x$ and smooth in $w$, such
that, on every occupied anchor cube,
\begin{equation}\label{test:mask-properties}
w_{i,\mathfrak t}=1\text{ on }G_{i,\mathfrak t},\qquad
\supp_w w_{i,\mathfrak t}\subset G_{i+1,\mathfrak t},\qquad
\|T_1\cdots T_jw_{i,\mathfrak t}\|_\infty\le C_jD_{l'}^j.
\end{equation}
Here the $T_j$ are any members of a fixed finite collection of smooth
tangent fields spanning the tangent spaces of the sphere.  The masks
are antipodally symmetric and constant as functions of $x$ on each
anchor cube; they are zero on unoccupied cubes and outside the root.
Moreover, the
$cD_{l'}^{-1}$-neighbourhood of $G_{i,\mathfrak t}$ lies in $G_{i+1,\mathfrak t}$, with
the same assertion for its closure after decreasing the fixed $c$.
\end{lemma}

\begin{proof}
Consider directions $w,w'$ at distance at most $cD_{l'}^{-1}$.
In a parallel test $[g,p]$, changing direction changes the distance
of a tested centre to the line by at most
\[
C_d2^{-g}|w-w'|
 \le C_dc2^{-p}R^{50\sigma}.
\]
Here $l'=p-g$; if $D_{l'}=1$, use
$2^{l'}\le R^{50\sigma}$.  In a projection test $[p,t]$ the change
in the tested scalar displacement is at most
\[
C_d2^{-p}|w-w'|
 \le C_dc2^{-t}R^{50\sigma},
\]
by the same argument with $l'=t-p$.  By contrast, the consecutive
width gap at the relevant fine depth $n$ is
\[
W_i(n)-W_{i+1}(n)=J_0^{-1}2^{-n}R^{500\sigma}.
\]
For fixed parameters and sufficiently large $N$, each perturbation
above is strictly less than half this gap.  Every configuration
counted at width $W_{i+1}$ in direction $w'$ was therefore counted at
width $W_i$ in direction $w$.  If $w$ is good at level $i$, then
$w'$ is good at level $i+1$.  The strict remaining margin also gives
the asserted inclusion for the closure of a smaller neighbourhood.

To construct the mask, take $r=c_0D_{l'}^{-1}$ with $c_0$ sufficiently
small, and the indicator of the spherical $2r$-neighbourhood of
$G_{i,\mathfrak t}$.  Convolve it on the sphere with a nonnegative rotationally
invariant smooth kernel of integral one, supported where the chordal
distance is less than $r$.  Such a kernel has angular derivative
$L^1$ bounds $C_jr^{-j}$, by rescaling in a fixed finite collection of
spherical charts.  The convolution equals one on $G_{i,\mathfrak t}$ and is
supported in the closure of its $3r$-neighbourhood, contained in
$G_{i+1,\mathfrak t}$.  The derivative bounds in
Equation~\eqref{test:mask-properties} follow.  Both tests and the
kernel are invariant under $w\mapsto-w$, so the mask is even.
The construction is made separately for each of the finitely many
anchor cubes; no spatial differentiability is required.
\end{proof}

\begin{lemma}[Products and retained supports]\label{test:product-symbols}
Consider a product with one factor for each occurrence in one or more
descriptor lists whose total size is at most $N^{C_*}$, with $C_*$ fixed
independently of $N$. A factor may be a mask $w_{i,\mathfrak t}$ for any
$0\le i<J_0$, or a derivative of order $q$ divided by $D_{l'}^q$, where
$l'$ is its nominal length. The indices may differ among factors. Assume
the total order of these prescribed derivatives is bounded independently
of $N$. A bounded number, independent of $N$, of measurable spatial
multipliers and smooth angular factors are also allowed, with fixed
numerical bounds for their sup norms and for all angular derivatives under
consideration. These bounds are common to the family of products in any
later supremum. Let $L$ be the largest nominal length, with $L=0$ for an
empty list, and set
\begin{equation}\label{test:total-symbol-scale}
D_*=R^\sigma\max(1,2^LR^{-50\sigma}).
\end{equation}
At each fixed angular derivative order $j$, the resulting product
$b$ satisfies
\[
\|T_1\cdots T_jb\|_\infty\le C_jD_*^j.
\]
For each factor with descriptor $\mathfrak t$ and index $i$, the support
of the product still imposes its $G_{i+1,\mathfrak t}$ condition, including
when that factor was already differentiated. All these assertions are
pointwise in the spatial variables. Constants may depend on the fixed
list-size exponent, prescribed derivative budget, and numerical bounds
for the extra factors, as well as on the previously fixed parameters.
\end{lemma}

\begin{proof}
An undifferentiated mask has norm at most one. Only a bounded number
of factors are prescribed derivatives, and their normalized derivatives
have bounds $C_jD_{l'}^j$ by Lemma~\ref{test:smooth-masks}, uniformly
over the allowed indices. At order $j$, the product rule gives
at most a polynomial in $N$ terms, by the assumed list-size bound
and because $j$ is fixed.  Each term is bounded by a
constant times $\max(1,2^LR^{-50\sigma})^j$.  For fixed
$\sigma>0$, its polynomial multiplicity is at most $R^{\sigma j}$
for sufficiently large $N$, giving
Equation~\eqref{test:total-symbol-scale}.  The order-zero bound has
no such multiplicity.  Finally, the support of a derivative of a
smooth function is contained in the support of the function.  Thus
none of the differentiations removes an imposed good-direction
condition.
\end{proof}

\subsection{Fixed lists under interval containment}

The individual tests and masks above are available for every descriptor.
We now choose a fixed list for each depth interval. The list for a larger
interval must contain the list for each subinterval, and each chosen
descriptor must have a selecting interval extending by its nominal length
on both sides of its anchor. We make each interval's own scale choice once
and reuse it whenever that interval contributes to a containing list.

Put $q_0=\lceil\kappa N\rceil$. Partition $\{0,\ldots,N\}$ into
consecutive blocks of $q_0$ indices, with the last block possibly
shorter. In each block choose, once and for all, a depth minimizing
$A$; break ties by a fixed rule. An interval $[a,t]$ is
\emph{terminal} if $t-a\le100q_0$. In a nonterminal interval select
$p$ by minimizing $A$ among the representatives of blocks wholly
contained in $[a+q_0,t-q_0]$, again using fixed tie breaking. Set
\[
l=\min(p-a,t-p).
\]
The choice is a \emph{prefix choice} if $2p\le a+t$, and a
\emph{suffix choice} otherwise.

For each interval $[a,t]$, form the list $\mathcal M(a,t)$ as follows.
For every nonterminal subinterval $[a',t']\subseteq[a,t]$, use its
fixed choice $p,l$ and include all parallel descriptors
\[
[g,p],\qquad p-l\le g<p.
\]
If the choice is a suffix choice, include also the projection
descriptor $[p,t']$. Repetitions within either type are removed.

For a membership $\mathfrak t\in\mathcal M(a,t)$, an \emph{origin}
means a nonterminal subinterval $[a',t']\subseteq[a,t]$ whose fixed
choice generated $\mathfrak t$. The origin witnesses this particular
membership; it is not part of the descriptor's identity, and different
memberships may use different origins. The initial entry will use one
additional parallel descriptor outside these lists. The individual test
and mask estimates apply to it without an origin.

\begin{lemma}[Descriptor geometry]\label{test:descriptor-properties}
For a nonterminal $[a,t]$, the union of the blocks eligible for its
choice is an integer interval $[a_*,t_*]$ with
\[
a+q_0\le a_*\le a+2q_0,\qquad
t-2q_0\le t_*\le t-q_0,
\qquad A_n\ge A_p\quad(a_*\le n\le t_*).
\]
The lists $\mathcal M(a,t)$ are nested under interval containment.
Across the scheduled lists, there are $O(1/\kappa)$ possible anchors
and polynomially many descriptors. For every
$\mathfrak t\in\mathcal M(a,t)$ with anchor $p$ and nominal length $l'$,
there is an origin $[a',t']\subseteq[a,t]$ such that
\begin{equation}\label{test:descriptor-origin}
[p-l',p+l']\subseteq[a',t'].
\end{equation}
That origin also supplies every parallel descriptor $[g,p]$ with
$p-l'\le g<p$. In particular $l'\le(t-a)/2$.
\end{lemma}

\begin{proof}
The eligible blocks are consecutive. The first starts less than one
block length after $a+q_0$, and the last ends less than one block
length before $t-q_0$. Nonterminality ensures that there are eligible
blocks. For an index in their union, its own block representative has
no larger profile value, and the selected $p$ has no larger value
than that representative. This proves the first assertions.

Each origin's choice is made independently of the interval in whose
list it is being included. Hence an origin contributing to a
smaller interval contributes the same descriptors to every larger
one, proving nesting. Anchors are block representatives, of which
there are at most $1+(N+1)/q_0$. A descriptor is specified by its
type, anchor, and one other integer depth. Their number is therefore
at most $2(N+1)(1+(N+1)/q_0)$.

A parallel descriptor coming from $[a',t']$ has $l'=p-g\le l$.
A projection descriptor occurs only for a suffix choice and then has
$l'=t'-p=l\le p-a'$. These observations prove
Equation~\eqref{test:descriptor-origin}; the list of parallel
descriptors supplied by the same origin contains all the asserted
ones. Deduplication preserves the conclusion for each membership:
by definition, at least one of the origins that generated a descriptor
in $\mathcal M(a,t)$ lies inside $[a,t]$. That origin may be chosen for
this membership independently of any choice made for another list.
\end{proof}

The scheduled lists satisfy the polynomial bound in
Lemma~\ref{test:product-symbols}, as do those lists with any fixed number
of extra descriptors. Their origin properties are needed only when using
the schedule; the individual test, exceptional-direction, and mask
estimates apply to every descriptor defined above.
